\subsubsection{Multiplier form of the action}
\label{p:sec:all:multiplier}

When the auxiliary fields cannot be taken on quotients, the construction has a variant in which nothing is inverted, the multiplier form of the action.
  It is written in terms of one field $\Sigma_n \in \Hn{n}$ per vertex, $\Sigma_1 := \Psi$, which fills all the slots of $\vertex{n}$, and one multiplier $\Theta_k \in \Hn{k}$ per such field, defined in \cref{p:app:conventions}:
  \begin{equation}
  \label{p:eq:all:multiplier-action}
  \caction[\mathrm{mult}]
    := \frac{1}{2} \, \pair{\Psi}{Q_B \Psi}
      + \sum_{n \ge 3} \frac{\gs^{n - 2}}{n!} \, \vertex{n}(\Sigma_n)
      + \sum_{k \ge 2} \pair[k]{\Theta_k}{\Sigma_k - \Sigma_{k - 1} \otimes \Psi},
\end{equation}
  which is cubic at all orders as it stands.
  The multipliers enforce the chain of constraints:
\begin{equation}
  \label{p:eq:all:constraints}
  \Sigma_k
    = \Sigma_{k - 1} \otimes \Psi,
  \qquad
  k \ge 2,
  \qquad
  \Sigma_1
    := \Psi,
\end{equation}
whose solution is $\Sigma_k = \Psi^k$ for every $k$.
The fields $\Sigma_n$ of this form are not the auxiliary fields of the construction above: they equal $\Psi^n$ exactly on the solution, whereas $\Sigma_k^* = \Psi^k + O(\gs)$, and they fill the elementary vertices rather than the reduced ones.
The multiplier equations are these constraints, and on their solution the multiplier terms vanish and the vertex terms are those of the original action, exactly and at every order, with no inverse of any vertex.
The equations of the auxiliary fields are linear in the multipliers and fix them from the top down, at any truncation of the tower or as a formal series in $\gs$.
The equation of $\Sigma_k$ involves $\Theta_k$, the multiplier $\Theta_{k + 1}$ contracted with $\Psi$ on one strand, and the vertex $\vertex{k}$ read as a $k$-string state through the pairing.
Each multiplier is thus determined by the multiplier of the next level, with real coefficients.
We have not worked out the gauge law and the Batalin--Vilkovisky structure of the multiplier form.

The multipliers pair with the auxiliary fields, so their ghost number is fixed by the saturation of the $k$-string pairing:
\begin{equation}
  \label{p:eq:all:ghost-multiplier}
  \Ngh(\Theta_k)
    = 3 k,
  \qquad
  \Ngh(\Sigma_k)
    = 2 k,
\end{equation}
one strand of ghost number three against each strand of ghost number two.
Their coefficients commute and their states have the parity $k \bmod 2$.
Their strands are read in the shifted parity of the output strands, as the pairing with a symmetric $\Sigma_k$ requires (\cref{p:app:conventions}).
With $\Theta_k^\star = \Theta_k$ the multiplier terms are real, and with the reality conditions~\eqref{p:eq:conv:reality-fields}, $\Psi^\star = - \Psi$ and $\Sigma_k^\star = (-1)^k \, \Sigma_k$, the whole action is real at all orders.

This form is the one to use when the auxiliary fields cannot be taken on quotients.
That is the case when the vertices do not annihilate the kernel of the quartic form, or of the nested operators at the higher levels, so that the conditions of the type~\eqref{p:eq:hs8:kernel-condition} fail.
It needs no kernel condition at all, since nothing is inverted.
Its price is that no auxiliary field has a quadratic term, nothing is eliminated by a Gaussian integral, and the exchange structure of the construction, the reduced vertices with it, is absent.
Between the two, the hybrid of \cref{sec:mult:hybrid} keeps the Hubbard--Stratonovich elimination on a complement of the kernel and a multiplier on the kernel directions only, and its reduced sextic vertex depends on the complement.
The datum of this form is then the complement together with the constraint sector and the reduced vertex.
A coupling of the multiplier to $(\Sigma_2, \Psi)$ at order $\gs^3$ can move the complement dependence from the vertex to the constraint, but not remove it.
This was followed through order $\gs^4$ in a finite-dimensional model, and transfers to the string vertices under the independence assumption on the exchange terms and the invertibility of the quartic form on the complement.
Beyond this order, we have not worked out the complement dependence of the later reduced vertices.
